\subsection{Canonical multiplicity purity, parameter weights and Lefschetz}
\label{sec:canonicalPureComparison}

The pure-point criterion applies directly to an individual
canonical coefficient, without assuming that the named right
operators are effective. The arithmetic assignment of the
source parameter must be specified; changing it can change
whether a raw coefficient admits a pure complex.

Work over \(\mathbb F_q\) with characteristic-zero
\(\ell\)-adic coefficients. Suppose a proposed canonical
product is \(T=Rf_*L\), where \(f\) is proper and \(L\)
is pure of integer complex weight \(W\). Let \(S_\gamma\)
be a geometrically absolutely simple perverse Weil sheaf,
pure of weight \(w_\gamma\), representing the proposed
output label. Its geometric isomorphism class is required
to be Frobenius stable, with its specified Weil form over
the base. A required finite extension of the base is made
explicitly; all later trace comparisons use that extension.

Define the geometric multiplicity spaces
\[
 V_i=\operatorname{Hom}_{\mathrm{Perv}(Y_{\overline{\mathbb F}_q})}
             (S_\gamma,{}^pH^i(T)),
 \qquad C_\gamma=\bigoplus_i V_i[-i],
\]
with conjugation Frobenius from the two specified Weil
structures. These are geometric Hom spaces, not their
Frobenius-fixed subspaces. The comparison being tested is
\begin{equation}\label{eq:canonicalMultiplicityTrace}
 f_\gamma(q^{n/2})
     =\sum_i(-1)^i\operatorname{Tr}(F^n\mid V_i)
\end{equation}
for every sufficiently large \(n\). No further inserted
endomorphism is part of this ordinary Euler trace.

\begin{finding}[\fid{F65}, the proper canonical comparison]
\label{f:properCanonicalMultiplicity}
Under these hypotheses, \(C_\gamma\) is a pure point
complex of weight \(W-w_\gamma\). If \(f_\gamma\) is an
integral Laurent polynomial, then
\[
                 (-1)^{W-w_\gamma}f_\gamma(-v)
                    \in\mathbb N[v,v^{-1}].
\]
In particular a nonzero coefficient divisible by
\(v-v^{-1}\) cannot satisfy this comparison. The same
conclusion holds after a monomial translation or an
overall sign, with the corresponding weight normalization.
\end{finding}

\begin{proof}
Proper pushforward preserves complex purity, and
\cite[Theorem~5.4.1, p.~141; Remark~5.4.9, p.~143]{BBD}
gives purity of \({}^pH^i(T)\) of weight \(W+i\).
Its geometric perverse sheaf is semisimple by
\cite[Theorem~5.3.8, p.~138]{BBD}. The canonical
\(S_\gamma\)-isotypic summand is
\(S_\gamma\otimes V_i\), with its conjugation Frobenius;
the evaluation isomorphism is Weil-equivariant. This is
the isotypic descent used explicitly in the proof of
\cite[Proposition~5.3.9, pp.~139--140]{BBD}.
It does not assert arithmetic semisimplicity: Frobenius
may have Jordan blocks on \(V_i\).

The summand is a perverse subquotient of a pure object,
so it is pure of weight \(W+i\)
\cite[Proposition~5.3.1, pp.~134--135]{BBD}.
On a smooth dense stratum \(U\) of the support, of
dimension \(d\), write
\(S_\gamma|_U=M[d]\), with \(M\) lisse and pure of
weight \(w_\gamma-d\). Regard these sheaves first on
their closed support. Closed pushforward and its inverse
on supported objects preserve both weight bounds; open
restriction has \(j^*=j^!\) and preserves both bounds
\cite[5.1.14, p.~128]{BBD}. Thus the lisse sheaf
\(M\otimes V_i\) has weight \(W+i-d\).
This does not use a false assertion about arbitrary
restrictions of pure complexes.

At a closed point of \(U\) of degree \(m\), let
\(\alpha\) be any Frobenius eigenvalue of \(M\), and
\(\beta\) any eigenvalue on \(V_i\). The product
\(\alpha\beta^m\) is an eigenvalue of the tensor
product. In every relevant embedding its modulus is
\(q^{m(W+i-d)/2}\), whereas that of \(\alpha\) is
\(q^{m(w_\gamma-d)/2}\). Hence
\[
                    |\beta|=q^{(W+i-w_\gamma)/2}.
\]
Thus \(V_i\) has weight \(W+i-w_\gamma\).
The shifted definition of purity makes
\(\bigoplus_iV_i[-i]\) pure of weight \(W-w_\gamma\).
The fully proved pure-point criterion of \fid{F57},
including its eventual-extension version, now gives
the asserted coefficient condition.

A nonzero Laurent polynomial divisible by \(v-v^{-1}\)
vanishes at \(-1\). Its asserted nonnegative transform
would vanish at \(1\), impossible for a nonzero
nonnegative Laurent polynomial. A monomial and an
overall sign preserve that vanishing; the corresponding
Tate and cohomological shifts change the weight in the
already specified way.
\end{proof}

\paragraph{The exact printed coefficient and a different parameter assignment.}
The individual canonical coefficient in VIII, Figure~2,
p.~34, and \(\S5.1.2\), p.~29, is
\begin{equation}\label{eq:VIIIrawCanonicalCoefficient}
 c(t)=-\frac{(t^2+t+1)(t-1)^2}{t^2}
           =-t^2+t+t^{-1}-t^{-2}.
\end{equation}
Here \(t\) denotes the source deformation parameter.
For \(t=q^n\) and \(v=q^{n/2}\), it is
\[
                    c(v^2)=-(v-v^{-1})^2(v^2+1+v^{-2}).
\]
It therefore fails the proper-pure multiplicity comparison
above. This conclusion concerns this individual coefficient;
it needs no effectivity premise for the named \(P_i\).
It does not refute VIII's signed/factored positivity.
The mixed nonproper model and proper normalized factor
were explicitly constructed in \fid{F62}.

If instead \(t=q^{n/2}\), the trace is \(c(v)\), and
\(-c(-v)=v^2+v+v^{-1}+v^{-2}\) is effective.
Let \(L_e\) be an arithmetic Weil line with eigenvalue
\(q^{e/2}\) and weight \(e\). Then
\begin{equation}\label{eq:alternatePureVIIIcoefficient}
 C_{\rm alt}=L_2[-1]\oplus L_1
                      \oplus L_{-1}[2]\oplus L_{-2}[3]
\end{equation}
is pure of weight one and has Euler trace \(c(v)\).
Its cohomological degrees are \(1,0,-2,-3\), respectively,
and the weights are \(2,1,-1,-2\).
The odd-weight lines are arithmetic coefficients; no
geometric integral half-Tate motive is asserted.
This is a pure point construction, not a canonical
proper-geometric realization. The two parameter rules
must not be silently identified.

\begin{finding}[\fid{F65}, the weight-assignment and Lefschetz boundary]
\label{f:canonicalLefschetzBoundary}
Assign \(t=q^{dn/2}\) for a nonzero integer \(d\).
The raw trace \(c(v^d)\) is the Euler trace of a pure
point complex exactly when \(d\) is odd; its possible
complex weights are then odd. For even \(d\), no
single pure weight is possible.

For every such nonzero \(d\), no pure point complex
with this trace carries a degree-two hard-Lefschetz
operator with a single center. The same obstruction
holds for \(\varepsilon v^m c(v^d)\), where
\(\varepsilon\in\{1,-1\}\) and \(m\in\mathbb Z\),
allowing the pure weight and the Lefschetz center to change.
Precisely, there are no such \(C,w,h\) and compatible
operator \(\ell\) for which
\[
 \ell^r:H^{h-r}(C)\xrightarrow{\sim}H^{h+r}(C)(r)
                       \qquad(r\geq0).
\]
\end{finding}

\begin{proof}
The polynomial \(c\) is invariant under inversion,
so replace \(d\) by \(|d|>0\). Its four roots in
the power-trace comparison correspond to exponents
\(-2d,-d,d,2d\). When \(d\) is even, all exponents
are even and the coefficients have mixed signs, which
fails the pure-point criterion for every weight.
When \(d\) is odd, \(-c((-v)^d)\) has four positive
coefficients, so precisely odd weights are possible.
The four-line construction places the root of exponent
\(e\) in degree \(e-w\), proving existence for each
such weight.

Suppose now that a pure complex has trace
\(\varepsilon v^m c(v^d)\). Finite power-sum uniqueness
fixes its roots. Purity puts each root in a unique
cohomological degree, so no additional roots can cancel
between opposite degrees. Consequently its only
nonzero cohomology groups, each one-dimensional, occur
in degrees
\[
 m-2d-w,\quad m-d-w,\quad m+d-w,\quad m+2d-w.
\]
If the required coefficient signs disagree with these
parities, purity already fails. Otherwise this support
is still forced. Hard Lefschetz makes the support
symmetric about \(h\). Its smallest and largest degrees
force \(h=m-w\). But \(H^h(C)=0\), while both
\(H^{h-2d}(C)\) and \(H^{h+2d}(C)\) are nonzero.
The asserted isomorphism \(\ell^{2d}\) between them
factors through \(H^h(C)\) after \(d\) applications
of \(\ell\), and is zero. This is impossible.
Jordan blocks cannot change the trace multiplicities
or this support argument.
\end{proof}

For a projective pushforward of a pure perverse input,
\cite[5.4.10, p.~144]{BBD} supplies the relative
hard-Lefschetz maps; taking a Frobenius-stable geometric
isotypic multiplicity preserves them. Thus the raw
coefficient also excludes that stronger projective
pure-perverse comparison for every nonzero integral
weight assignment of the source variable. General
proper pure complexes do not automatically satisfy
this additional Lefschetz premise. Nonproper mixed
operations, inserted traces, and the normalized
positive factor remain distinct possibilities.

% Separate paragraph for Part III.
\paragraph{Canonical multiplicities and parameter weights in Grothendieck's language.}
The spaces \(V_i\) are actual Weil representations.
Their finite coefficient complex and Euler class use
SGA~6, I, Definition~4.7, pp.~121--122, and
Example~4.8, p.~122. The split-torus effective cone
uses SGA~3, I, 4.7.2--3 and cote~90, scan~17;
SGA~6, V, \S\S2--3 and cote~133, scan~18 distinguish
actual symmetric coefficient operations and virtual
classes \cite{SGA3,SGA6,Archive}.
No commutative \(\lambda\)-structure on the convolution
algebra is inferred.

The proper pushforward, pure perverse subquotients and
geometric isotypic multiplicities use BBD~5.1.14,
p.~128; 5.3.1, pp.~134--135; 5.3.8, p.~138;
the proof of 5.3.9, pp.~139--140; 5.4.1, p.~141;
and 5.4.9--5.4.10, pp.~143--144 \cite{BBD}.
The Jordan example after 5.3.8 distinguishes geometric
from arithmetic semisimplicity. SGA~5, III, 4.7
provides the trace framework \cite{SGA5}; it does
not identify a canonical basis. The new deductions
are the precisely normalized canonical-coefficient
obstruction, the alternate pure point model, and the
missing-middle Lefschetz test. They are not archival
claims or an assertion that VIII imposes the tested
geometric comparison.
